\documentclass[11pt,a4paper]{article}

\usepackage[margin=2.5cm]{geometry}
\usepackage{hyperref}
\usepackage{listings}
\usepackage{xcolor}

\lstset{
  basicstyle=\ttfamily\small,
  backgroundcolor=\color{gray!10},
  frame=single,
  breaklines=true,
  columns=fullflexible,
}

\title{\textbf{CAM-Expert 1.2 --- Build Information}}
\author{CAM-Expert Project}
\date{}

\begin{document}

\maketitle
\tableofcontents
\newpage

% ---------------------------------------------------------------------------
\section{Requirements}

\begin{itemize}
  \item \textbf{OS:} Windows 10 or 11 (primary runtime target)
  \item \textbf{CMake:} 3.20 or newer
  \item \textbf{Compiler:} C++20-capable toolset
    \begin{itemize}
      \item Visual Studio 2019 (toolset v142), 2022 (v143), or 2026 (v180)
      \item MinGW-w64 on Windows (GCC or Clang with Win32 headers)
    \end{itemize}
  \item \textbf{System libs:} Win32 platform libraries supplied by the Windows SDK
    (\texttt{opengl32}, \texttt{glu32}, \texttt{gdi32}, \texttt{user32},
     \texttt{comctl32}, \texttt{comdlg32}, \texttt{shell32}, \texttt{ole32},
     \texttt{oleaut32}, \texttt{uuid}, \texttt{dwmapi}, \texttt{uxtheme})
  \item \textbf{Optional:} OpenMP --- detected automatically; project builds without it
    if unavailable (\texttt{CAMEXPERT\_USE\_OPENMP} is only defined when found)
\end{itemize}

% ---------------------------------------------------------------------------
\section{Configure and Build (Windows)}

\subsection{Visual Studio 2026}
\begin{lstlisting}[language=bash]
cmake -B build -G "Visual Studio 18 2026" -A x64
cmake --build build --config Release
\end{lstlisting}

\subsection{Visual Studio 2022}
\begin{lstlisting}[language=bash]
cmake -B build -G "Visual Studio 17 2022" -A x64
cmake --build build --config Release
\end{lstlisting}

\subsection{Visual Studio 2019}
\begin{lstlisting}[language=bash]
cmake -B build -G "Visual Studio 16 2019" -A x64
cmake --build build --config Release
\end{lstlisting}

\subsection{MinGW-w64 on Windows}
\begin{lstlisting}[language=bash]
cmake -B build -G "MinGW Makefiles" -DCMAKE_BUILD_TYPE=Release
cmake --build build
\end{lstlisting}

% ---------------------------------------------------------------------------
\section{CMake Presets}

Presets are defined in \texttt{CMakePresets.json} and are host-conditioned to
Windows (\texttt{hostSystemName == Windows}).

\begin{lstlisting}[language=bash]
cmake --list-presets

# Configure
cmake --preset windows-vs2026   # Visual Studio 2026
cmake --preset windows-vs2022   # Visual Studio 2022
cmake --preset windows-vs2019   # Visual Studio 2019
cmake --preset windows-mingw-release  # MinGW-w64

# Build
cmake --build --preset release-vs2026
cmake --build --preset release-vs2022
cmake --build --preset release-vs2019
cmake --build --preset release-mingw
\end{lstlisting}

\noindent On non-Windows hosts the preset conditions will not match; use the
explicit \texttt{-G} generator commands from Section~2 instead.

% ---------------------------------------------------------------------------
\section{Linux / Container Notes}

A \texttt{Dockerfile} is provided for containerised configure/build
reproducibility:

\begin{lstlisting}[language=bash]
docker build -t camexpert-ci .
docker run --rm camexpert-ci
\end{lstlisting}

\noindent \textbf{Important:} CAM-Expert is a Win32 application that requires
\texttt{windows.h} and the Windows SDK headers.  A Linux-native compile is
expected to fail at the \texttt{windows.h} include unless a
Windows-targeting cross-compilation toolchain is configured.  The Docker
container is therefore useful only for CMake configure-step smoke checks,
not for producing a runnable binary.

% ---------------------------------------------------------------------------
\section{CMake Build Variables}

\begin{center}
\begin{tabular}{lll}
\hline
\textbf{Variable} & \textbf{Default} & \textbf{Description} \\
\hline
\texttt{CMAKE\_BUILD\_TYPE} & (generator-dependent) & \texttt{Release} / \texttt{Debug} for Make-based generators \\
\texttt{CAMEXPERT\_USE\_OPENMP} & auto-detected & Defined to \texttt{1} when OpenMP is found \\
\hline
\end{tabular}
\end{center}

% ---------------------------------------------------------------------------
\section{Build Definitions (Always Active on Windows)}

\begin{itemize}
  \item \texttt{UNICODE}, \texttt{\_UNICODE} --- enable wide-character Win32 API
  \item \texttt{WIN32\_LEAN\_AND\_MEAN} --- reduces Windows header bloat
  \item \texttt{NOMINMAX} --- prevents \texttt{min}/\texttt{max} macro conflicts
  \item \texttt{\_CRT\_SECURE\_NO\_WARNINGS} --- suppresses MSVC CRT deprecation warnings
\end{itemize}

% ---------------------------------------------------------------------------
\section{Compiler Options}

\begin{center}
\begin{tabular}{ll}
\hline
\textbf{Toolset} & \textbf{Flags} \\
\hline
MSVC & \texttt{/W4 /MP} \\
GCC / Clang & \texttt{-Wall -Wextra -Wpedantic} \\
\hline
\end{tabular}
\end{center}

% ---------------------------------------------------------------------------
\section{Install}

\begin{lstlisting}[language=bash]
cmake --install build --prefix <destination>
\end{lstlisting}

\noindent The executable \texttt{CAMExpert.exe} is installed to
\texttt{<destination>/bin/}.

% ---------------------------------------------------------------------------
\section{Validation and Smoke Testing}

There is currently \textbf{no automated unit-test suite}.  The recommended
post-build validation workflow is:

\begin{enumerate}
  \item CMake configure smoke test (no errors during configure).
  \item Build smoke test on a Windows toolchain (successful compile and link).
  \item Manual UI and toolpath smoke run (wireframe / surface / solid actions).
  \item CAM generation checks (2D, 3D, and multi-axis as applicable).
  \item Post-output sanity check (inspect generated NC file contents).
  \item Verify / backplot / machine-simulation smoke checks.
\end{enumerate}

% ---------------------------------------------------------------------------
\section{Version Synchronisation}

During release hardening, the version number must be kept consistent across:

\begin{itemize}
  \item \texttt{CMakeLists.txt} --- \texttt{project(CAMExpert VERSION 1.2.0 ...)}
  \item \texttt{src/Application.h} --- application version string constants
  \item \texttt{src/resources/CAMExpert.rc} --- Windows version resource fields
\end{itemize}

% ---------------------------------------------------------------------------
\section{Troubleshooting}

\begin{description}
  \item[\texttt{windows.h: No such file or directory}]
    You are compiling outside a Windows-native or Windows-cross toolchain
    environment.  This is expected on Linux unless Win32 headers are available.

  \item[Visual Studio toolset mismatch (\texttt{MSB8020})]
    Use the CMake generator that matches the Visual Studio version you have
    installed, or add the required toolset component via Visual Studio
    Installer.

  \item[CMake preset not available]
    Presets are host-conditioned for Windows.  On non-Windows hosts, use the
    explicit \texttt{-G} generator commands from Section~2.

  \item[OpenMP not found]
    OpenMP is optional.  The project builds and runs without it; only certain
    3D toolpath projection loops lose parallel acceleration.
\end{description}

\end{document}
